\section{Explicit spectral bounds and proof}
\label{app:note-spectrum}

The main theorem concerns the actual finite tanh dictionary. An auxiliary
center integral exposes its gamma dependence; lattice and exterior-center
allowances transfer its spectrum to that dictionary. This appendix specifies
each approximation and the resulting endpoints.

\paragraph{Finite setup.}
Let $z_1,\ldots,z_M$ be distinct training points with $M\ge2$, $W$ consecutive centers on
$c_0+h\mathbb Z$ cover their interval, and $\gamma>0$. Define
\begin{equation}
 \Phi_\gamma=[\mathbf1,(\tanh(\gamma(z_a-x_j)))_{a,j}],\qquad
 K_\gamma=\frac1M\Phi_\gamma\Phi_\gamma^\top,\qquad
 y=(f(z_a))_{a=1}^M\ne0.
\end{equation}
Readout GD minimizes $\|\Phi_\gamma w-y\|^2/(2M)$.
The theorem permits arbitrary distinct samples; the reported experiments use
uniform samples. Write $v_0=\mathbf1/\sqrt M$ and let $Q$ have orthonormal
columns spanning $v_0^\perp$. This projection removes the common constant
sample pattern, not any target assumption. Put
$t_\gamma(c)=(\tanh(\gamma(z_a-c)))_a$.

\paragraph{The gamma-dependent correlation model.}
Its center-integrated, mean-zero kernel is
\begin{equation}
 H_\gamma=\frac1{hM}\int_{\mathbb R}
 Q^\top t_\gamma(c)t_\gamma(c)^\top Q\,dc
 =-\frac2{hM}Q^\top[d_{ab}\coth(\gamma d_{ab})]_{a,b}Q,
 \quad d_{ab}=z_a-z_b.
 \label{eq:note-H}
\end{equation}
The diagonal of the bracketed matrix is $1/\gamma$. Choose a finite set
$\mathcal F$ of exterior lattice centers absent from the actual dictionary,
leaving two unretained tails starting at $c_L<\min z_a$ and
$c_R>\max z_a$. Subtract their retained contribution:
\begin{equation}
 T_\gamma=\frac1M\sum_{c\in\mathcal F}
 Q^\top t_\gamma(c)t_\gamma(c)^\top Q,\qquad
 S_\gamma=H_\gamma-T_\gamma.
 \label{eq:note-S}
\end{equation}
Let $\beta_1\ge\cdots\ge\beta_{M-1}$ be the eigenvalues of $S_\gamma$.
For any $0<\theta<\pi/2$, the lattice and remaining-tail allowances are
\begin{align}
 \delta_\gamma&=
 \frac{8\gamma\|z-\bar z\mathbf1\|^2\sec^4\theta}
 {3hM[\exp(2\pi\theta/(\gamma h))-1]},\qquad
 \bar z=M^{-1}\sum_a z_a,\label{eq:note-delta}\\
 \tau_\gamma&=
 \frac{4[e^{-4\gamma(c_R-z_{\max})}
       +e^{-4\gamma(z_{\min}-c_L)}]}{1-e^{-4\gamma h}}.
 \label{eq:note-tau}
\end{align}
They give the matrix enclosure
\begin{equation}
 S_\gamma-(\delta_\gamma+\tau_\gamma)I
 \preceq Q^\top K_\gamma Q\preceq S_\gamma+\delta_\gamma I.
 \label{eq:note-matrix-enclosure}
\end{equation}

\paragraph{Relative-eigenvalue endpoints.}
Normalize using scalar bounds $\ell_\gamma\le\mu_1\le L_\gamma$.
The following require only matrix-vector products with the finite dictionary:
\begin{align}
 \ell_\gamma&=v_0^\top K_\gamma v_0\ge1,
 &b_\gamma&=\|Q^\top K_\gamma v_0\|,
 &c_\gamma&=\beta_1+\delta_\gamma,\\
 L_\gamma&=\frac{\ell_\gamma+c_\gamma+
 \sqrt{(\ell_\gamma-c_\gamma)^2+4b_\gamma^2}}2.
 \label{eq:note-normalization}
\end{align}
Set $\underline\rho_1=\overline\rho_1=1$ and $\beta_M=0$.
The endpoints in Theorem~\ref{thm:note-spectrum} are
\begin{equation}
 \boxed{\underline\rho_i=
 \frac{[\beta_i-\delta_\gamma-\tau_\gamma]_+}{L_\gamma},
 \qquad
 \overline\rho_i=\min\!\left\{1,
 \frac{\beta_{i-1}+\delta_\gamma}{\ell_\gamma}\right\},
 \quad 2\le i\le M.}
 \label{eq:note-endpoints}
\end{equation}
Exact zeros known from $\operatorname{rank}K_\gamma\le W+1$ may have both
endpoints set to zero. No target or optimization trajectory enters these
spectral bounds.

\subsection{Derivation of the explicit gamma dependence}
The identity
$1-\tanh u\tanh v=\coth(u-v)(\tanh u-\tanh v)$ gives
\begin{equation}
 \int_{\mathbb R}[1-\tanh(\gamma(z_a-c))\tanh(\gamma(z_b-c))]dc
 =2d_{ab}\coth(\gamma d_{ab}).
\end{equation}
Multiplication by $Q^\top$ and $Q$ removes the constant term and proves
\eqref{eq:note-H}. The integral is finite after projection because
$Q^\top t_\gamma(c)$ decays at both ends.

The smoothing interpretation can also be made exact. With
$s_\gamma(x)=\frac\gamma2\sech^2(\gamma x)$,
$\tanh(\gamma\,\cdot)=s_\gamma*\sign$. Its Fourier multiplier is
\begin{equation}
 A_\gamma(\omega)=\widehat s_\gamma(\omega)
 =\frac{\pi|\omega|/(2\gamma)}{\sinh(\pi|\omega|/(2\gamma))},
 \qquad A_\gamma(0)=1.
 \label{eq:note-attenuation}
\end{equation}
For completeness, substitution $r=e^{2\gamma x}$ in the Fourier transform
gives $\int_0^\infty r^{-i\omega/(2\gamma)}(1+r)^{-2}dr$.
The beta integral evaluates this as
$\Gamma(1-i\omega/(2\gamma))\Gamma(1+i\omega/(2\gamma))$,
which equals \eqref{eq:note-attenuation} by the reflection formula.
For $u=Qv$, the step combination $\sum_a u_a\sign(c-z_a)$ has compact
support, and its derivative is $2\sum_a u_a\delta_{z_a}$. Parseval gives
\begin{equation}
 v^\top H_\gamma v=\frac2{\pi hM}\int_{\mathbb R}
 \frac{A_\gamma(\omega)^2}{\omega^2}
 \left|\sum_a(Qv)_a e^{-i\omega z_a}\right|^2d\omega.
 \label{eq:note-directional-filter}
\end{equation}
The apparent zero-frequency singularity is removable since $\sum_a(Qv)_a=0$.
Thus gamma acts on a direction's frequency content through the explicit
weight $A_\gamma^2$. Since $t/\sinh t$ decreases for $t>0$, increasing gamma
increases every directional response of $H_\gamma$ and hence its ordered
eigenvalues. This order concerns the integral kernel; exterior subtraction
and normalization need not preserve monotonicity of every finite ratio.

\subsection{Transfer to the finite dictionary}
For $u\perp\mathbf1$, the function
$g(\zeta)=\sum_a u_a\tanh(\gamma(\zeta-z_a))$ is analytic for
$|\operatorname{Im}\zeta|<\pi/(2\gamma)$.
Subtracting the common translate at $\bar z$ and integrating derivatives,
Minkowski and Cauchy--Schwarz yield
\begin{equation}
 \int_{\mathbb R}|g(t\pm i\theta/\gamma)|^2dt
 \le\frac{4\gamma}{3}\sec^4\theta
 \|z-\bar z\mathbf1\|^2\|u\|^2.
 \label{eq:note-strip}
\end{equation}
Here $|\sech^2(t+i\theta)|\le\sec^2\theta\sech^2t$ and
$\int\sech^4t\,dt=4/3$. Shift the contour for the analytic function
$g(\zeta)^2$, not $|g(\zeta)|^2$: its transform is bounded by the
right side of \eqref{eq:note-strip} times $e^{-\theta|\omega|/\gamma}$.
Poisson summation at frequencies $2\pi k/h$ then gives
\begin{equation}
 \left\|\frac1M\sum_{c\in c_0+h\mathbb Z}
 Q^\top t_\gamma(c)t_\gamma(c)^\top Q-H_\gamma\right\|
 \le\delta_\gamma.
\end{equation}
The geometric sum over $k\ne0$ produces the denominator in
\eqref{eq:note-delta}.

Deleting exterior centers converts the infinite lattice to the actual
dictionary. The retained deletions give $T_\gamma$. For a remaining center
$c\ge c_R$, projection removes the constant limit and
$1-\tanh t\le2e^{-2t}$ gives an outer-product norm, divided by $M$, of at
most $4e^{-4\gamma(c-z_{\max})}$.
Summing this geometric series and its left-tail counterpart bounds the
remaining positive deletion by $\tau_\gamma I$. This proves
\eqref{eq:note-matrix-enclosure}.

If $\kappa_i$ are the eigenvalues of $Q^\top K_\gamma Q$, min--max and
codimension-one interlacing give
\begin{equation}
 \beta_i-\delta_\gamma-\tau_\gamma\le\kappa_i
 \le\beta_i+\delta_\gamma,
 \qquad \mu_i\ge\kappa_i\ge\mu_{i+1}.
\end{equation}
The Rayleigh quotient gives $\mu_1\ge\ell_\gamma$.
In the decomposition $\operatorname{span}(v_0)\oplus v_0^\perp$,
the diagonal blocks of $K_\gamma$ are bounded above by
$\ell_\gamma$ and $c_\gamma I$, and the off-diagonal block has norm
$b_\gamma$. Its quadratic form is therefore bounded by the largest
eigenvalue of $\left[\begin{smallmatrix}\ell_\gamma&b_\gamma\\
b_\gamma&c_\gamma\end{smallmatrix}\right]$, namely $L_\gamma$.
Dividing the interlacing bounds by these normalization endpoints proves
\eqref{eq:note-endpoints} and the main theorem.

\subsection{Classical output-error consequence}
For $r_n=\Phi_\gamma w_n-y$, normalized least-squares GD gives
$r_{n+1}=(I-\eta K_\gamma)r_n$ and $r_0=-y$. Expansion in the actual
finite-kernel eigenvectors gives
\begin{equation}
 e_n^2=\sum_i p_i(1-\eta\mu_i)^{2n}
 =\sum_i p_i(1-\rho_i/2)^{2n}
 \ge\sum_i p_i(1-\overline\rho_i/2)^{2n}.
\end{equation}
The last inequality follows because the factor decreases with
$\rho_i\in[0,1]$. This proves Corollary~\ref{cor:note-output}.
Our step is a conservative fraction of the stability limit $2/\mu_1$;
any nonoscillatory step $0<\eta\mu_1\le1$ instead uses $\eta\mu_1\rho_i$.
Only exact nullspace energy gives a permanent floor. For $\mu_i>0$,
$u_i=\Phi_\gamma(\Phi_\gamma^\top u_i/(M\mu_i))$ is representable,
however small its rate. The target weights require the actual eigenvectors;
the theorem bounds eigenvalues, not their target alignment.

\section{Evaluation, accuracy, and empirical scope}
\label{app:note-evaluation}

\paragraph{Which quantities are computed.}
The theorem endpoints use the explicit center integral, exterior-center
subtraction, and scalar normalization bounds. A separate SVD of the actual
finite feature matrix supplies the validation eigenvalues and target weights
$p_i$. The plotted lower error uses the theorem's upper rates, not actual
rates. The spectral figure selects three ranks for readability; the error
bound includes every numerically resolved target projection.

The integral is evaluated through a positive quadrature factor rather than
by subtracting nearly equal entries of \eqref{eq:note-H}. There is no Fourier
truncation. Truncating the center integral at
$\pm(\max_a|z_a|+P/\gamma)$ omits a positive matrix of norm at most
$2e^{-4P}/(h\gamma)$, added to the numerical upper endpoint.
The exterior remainder enters the lower endpoint. Compressing the combined
integral and exterior factors $Z$ to rank $r$ introduces the two-sided bound
$2\sigma_1(Z)\sigma_{r+1}(Z)+\sigma_{r+1}(Z)^2$ on their signed Gram matrix.
Quadrature refinement and an empirical arithmetic guard are checked
separately; the figures are FP64 evaluations, not interval-arithmetic
certificates. Refining quadrature order from 10 to 16 and padding from
20 to 24 changes the lower error curves by at most $9.6\times10^{-11}$.

The remaining analytic slack comes from the interlacing index shift, the
normalization bounds $\ell_\gamma,L_\gamma$, and additive allowances.
Small absolute errors can dominate extremely small eigenvalues; the
enclosure is valid but need not be informative for every geometry or mode.
In this experiment it is informative both for the plotted ratios and for
the full error trajectory. Omitting unresolved target energy from the lower
curve only weakens it; it is not interpreted as an exact approximation floor.

\paragraph{Target, geometry, and measurements.}
All displayed data use
\begin{equation}
 f(x)=\frac{\sin(2\pi x)+\frac12\sin(6\pi x)+\frac14\sin(14\pi x)}
 {\sqrt{21/32}},\qquad x\in[-1,1].
\end{equation}
The total hidden width is $W=512$, comprising $N=467$ interior cells and
$R=22$ halo centers on each side, so $h=2/467$ and $W=N+2R+1$.
There are 2,048 uniform training midpoints, 4,096 validation points, and
8,192 dense evaluation points. Every reported error is relative output
$L^2$ error on its stated grid, not squared loss or an EMA.
The dense grid checks resolution; it is not an untouched generalization test.
Frozen GD starts its readout at zero and executes five million full-batch
FP64 updates with $\eta=1/(2\mu_1)$ at each bandwidth.
Figure~\ref{fig:note-training}A displays training error. Its spectral
reference agrees with executed GD within $2.9\times10^{-14}$.

Joint training optimizes all slopes, centers (through affine biases), and
readout parameters for five million updates on five paired seeds. One
recipe per optimizer is selected by median final validation error over all
five seeds. The selected Adam and GD schedules are full-horizon cosine
decays with initial learning rates 0.05 and 0.2, respectively.
Frozen Adam uses the same final-validation selection criterion.
Panel B shows raw training errors summarized for display; quoted joint and
frozen-Adam endpoint values use the dense grid. Display bins retain extrema,
including spikes, rather than replacing the comparison with first tolerance
crossings. Panel C measures $h|a_j|$ for learned features
$\tanh(a_jx+b_j)$, using the original reference spacing $h$.

\paragraph{What the joint experiment establishes.}
Adam's median scaled RMS slope is $0.0582$, and its 99th percentile is
$0.2466$; the five seeds have $6,5,5,4,7$ features at or above $1/4$.
These statistics describe heterogeneous geometry, not an effective uniform
bandwidth. The uniform $1/4$ dictionary is an empirical comparison, not a
necessary scale for every target or learned dictionary.

Joint Adam continues improving late in training; its median endpoint error
falls about 28\% between the two- and five-million-update experiments.
The result is a precision gap under a large finite budget, not permanent
failure. Checked direct fits and an additional two-million-update readout
restart on the learned five-million-update features barely improve their
error. Multiplying the learned slopes by 4 or 16 at fixed centers worsens
the median refitted error to 0.0916 or 0.139, respectively. These interventions
change nonnested feature spaces: they support investigating slopes and
centers together, not a theorem that increasing slopes always helps.
Neither the fixed-kernel proof nor its scalar GD rates apply to evolving
features or Adam's adaptive preconditioner. The next theoretical question
is how evolving features acquire a useful geometry.

\paragraph{Reproducibility.}
The source and numerical provenance for the new spectral figure are
\path{experiments/expD36_frozen_gamma_probe/section34_spectrum_figure.py}
and \path{output/diagnostics/section34_long_horizon/main_5m/spectrum/relative_rates.json}.
The existing three-panel figure, full protocols, checks, and evidence links
are indexed in \path{output/diagnostics/section34_long_horizon/README.md}.
Both figures reuse completed runs; preparing this note launches no new
training experiments.
